% ==============================================================================
% PLANTILLA MAESTRA GENERAL - UCNL
% ==============================================================================
% Uso:
% 1. Duplica este archivo y renombra la copia por materia, semana y actividad.
% 2. Actualiza documenttitle, documentsubtitle, documentsubject, coursename,
%    coursecode, universitydepartment y datos de la figura docente.
% 3. Lee la planeacion semanal y NO pegues las instrucciones completas dentro del
%    producto final. Usa la planeacion como lista de cumplimiento.
% 4. Revisa la bibliografia indicada, toma notas, registra citas y redacta con
%    comprension propia.
% 5. Identifica el producto academico solicitado y construyelo dentro del
%    documento siempre que sea posible.
% 6. Entrega en PDF con la nomenclatura indicada en el aula.
%
% Formato institucional recomendado:
% - Fuente tipo Helvetica, equivalente visual a Arial.
% - Interlineado 1.5.
% - Citas autor-anio con natbib: usar \citep{clave} o \citet{clave}.
% - Referencias en bibliografia BibTeX, formato APA 7 en cuerpo y referencias.
% - Caratula, resumen breve, desarrollo, producto solicitado, conclusion y
%   bibliografia.
% - Redaccion academica clara, coherente, cohesionada, formal y sin faltas.
% - La portada puede llevar una marca de agua institucional discreta; ajustar
%   las variables coverwatermark... si se requiere apagarla o cambiar intensidad.
%
% Criterios generales de cumplimiento:
% - Responder exactamente a la actividad solicitada en la planeacion.
% - Identificar conceptos, elementos, criterios o categorias de forma precisa.
% - Analizar, interpretar y tomar postura; no solo copiar fuentes.
% - Organizar la informacion por jerarquia, secuencia y criterios claros.
% - Incluir conclusion personal cuando la actividad lo permita o lo solicite.
% - Usar citas y referencias en APA 7.
% - Evitar copiar las indicaciones de la actividad en el cuerpo final.
%
% Regla para productos visuales:
% - Si la actividad pide tabla, cuadro comparativo, esquema, mapa conceptual,
%   mapa mental, linea de tiempo, diagrama, matriz, lista de verificacion, red
%   conceptual, SQA, organizador grafico o producto similar, realizarlo
%   preferentemente con LaTeX/TikZ para mantener nitidez, evidencia editable y
%   diseno controlado.
% - Si la consigna exige Canva, Miro u otra herramienta externa, insertar una
%   captura nitida y conservar la explicacion dentro del documento.
% - Para tablas extensas, usar pagina limpia, sin encabezados ni pies, landscape,
%   letra pequena y restaurar el estilo al terminar.
%
% Criterios de redaccion personal:
% - No escribir para llenar espacio; escribir para sostener una idea.
% - Convertir informacion en criterio: que significa, como funciona, que problema
%   resuelve, que limite presenta y que consecuencia produce.
% - Estructura mental recomendada: problema -> sistema -> consecuencia -> postura.
% - Usar pensamiento de diagnostico: entrada -> proceso -> salida.
% - Cada parrafo debe tener funcion: presentar, explicar, comparar, valorar o
%   concluir.
% - Usar primera persona con moderacion academica: "Considero que",
%   "Desde esta perspectiva", "La posicion que sostengo es". Evitar "yo creo".
%
% Notas de enfoque personal segun enfoque_personal.txt:
% - La voz debe ser academica, tecnica, directa, estrategica y funcional.
% - El texto no debe sonar como estudiante pasivo que recopila: debe sonar como
%   alguien que interpreta, ordena, cruza datos y busca consecuencia practica.
% - Antes de redactar, identificar la funcion de cada seccion: abrir problema,
%   definir criterio, demostrar con evidencia, comparar, valorar o concluir.
% - Evitar frases genericas como "es importante porque si" o "desde siempre".
%   Sustituirlas por relaciones verificables: causa, efecto, limite, costo,
%   funcion, riesgo, implicacion o mejora.
% - Cada concepto debe pasar de definicion a analisis. Formula util:
%   concepto -> funcion -> riesgo si se aplica mal -> consecuencia -> postura.
% - La postura personal no debe ser opinion suelta. Debe tener tres piezas:
%   posicion, razon y efecto. Ejemplo: "Considero que X, porque Y; por ello Z".
% - Integrar experiencia o criterio propio sin perder formalidad: no narrar la
%   vida personal, sino usar mirada practica para explicar por que el tema
%   importa en educacion, derecho, instituciones, comunicacion o profesion.
%
% Estilo observado en actividades ya realizadas:
% - Abrir con una tesis concreta:
%   "El Derecho no solo ordena conductas: organiza poder, define limites y
%   distribuye consecuencias."
% - Formular el problema en negativo y positivo:
%   "El problema no es solo..., sino..."
% - Convertir el producto visual en argumento:
%   "El cuadro no solo recopila datos; reconstruye la funcion, el problema y el
%   punto que requiere correccion."
% - Cerrar con aprendizaje transferible:
%   "No basta con saber que dice la norma; tambien debe analizarse a quien
%   protege, a quien excluye y que razones permiten defenderla publicamente."
%
% Nivel cognitivo esperado:
% - Recordar: usar solo para definir conceptos basicos. No debe ser el nivel
%   dominante en una actividad ponderable.
% - Comprender: explicar con palabras propias y relacionar con el tema.
% - Aplicar: llevar el concepto a un caso, norma, texto, correo, dilema o hecho.
% - Analizar: separar elementos, detectar causas, efectos, tensiones y limites.
% - Evaluar: sostener una postura con criterios, fuentes y consecuencias.
% - Crear: elaborar producto propio: mapa, cuadro, linea de tiempo, ensayo,
%   propuesta, reformulacion, glosario, portafolio o solucion argumentada.
%
% Regla de nivel minimo:
% - Si la actividad vale puntos y pide producto, no quedarse en recordar o
%   comprender. Debe llegar por lo menos a aplicar y analizar; si pide conclusion
%   o postura, debe llegar a evaluar.
%
% Uso de las 100 tecnicas didacticas:
% - Esta plantilla conserva el manual "Tecnicas didacticas del aprendizaje" que
%   se incorporo en las plantillas de Filosofia del Derecho, Etica y Moral
%   juridica, y Redaccion en contextos virtuales.
% - Cuando la planeacion mencione una tecnica de la coleccion UnADM, identificar:
%   definicion de la tecnica, producto esperado, pasos de construccion, criterios
%   de revision y cierre reflexivo.
% - La tecnica didactica no debe verse como adorno. Debe mostrar seleccion,
%   jerarquia, relacion entre ideas y lectura propia.
% - Siempre que sea viable, construir la tecnica en LaTeX/TikZ: cuadro
%   comparativo, mapa conceptual, esquema, mapa mental, linea de tiempo, matriz,
%   red conceptual, diagrama de flujo, tabla SQA o lista de verificacion.
%
% Checklist antes de entregar:
% [ ] El titulo, materia, semana, docente y fecha estan actualizados.
% [ ] El objetivo coincide con la planeacion.
% [ ] El producto solicitado aparece visible, rotulado y completo.
% [ ] La introduccion plantea un problema o postura, no solo anuncia el tema.
% [ ] El desarrollo explica, conecta y analiza.
% [ ] Hay citas APA autor-anio dentro del texto.
% [ ] La bibliografia final coincide con las fuentes citadas.
% [ ] La conclusion responde que se comprendio y que postura se sostiene.
% [ ] El nivel cognitivo llega al menos a aplicacion y analisis.
% [ ] Si hay conclusion/postura, incluye posicion, razon y consecuencia.
% [ ] La tecnica didactica solicitada esta construida y explicada.
% [ ] No aparecen instrucciones copiadas de la planeacion dentro del producto.
% [ ] Ortografia, acentos, sintaxis y nombres propios revisados.
% [ ] Si se uso IA como apoyo, el texto fue comprendido, revisado y apropiado.
% ==============================================================================

% CREACION DEL DOCUMENTO
\documentclass[
	spanish,
	letterpaper, oneside
]{article}

% ==============================================================================
% INFORMACION DEL DOCUMENTO
% Actualizar en cada actividad.
% Ejemplos:
% \def\documenttitle {Cuadro comparativo de conceptos fundamentales}
% \def\documentsubtitle {Actividad X - Contabilidad I}
% \def\documentsubject {Informe academico UCNL}
% ==============================================================================
\def\documenttitle {Actividad 2 - Contabilidad I}
\def\documentsubtitle {Cuestionario sobre ciclo contable y estados financieros}
\def\documentsubject {Actividad 2}

\def\documentauthor {Martin Jonathan de la Cruz}
\def\coursename {Contabilidad I}
\def\coursecode {CON-I}

\def\universityname {Universidad Ciudadana de Nuevo Leon}
\def\universityfaculty {Programa academico UCNL}
\def\universitydepartment {Contabilidad I}
\def\universitydepartmentimage {departamentos/logo-ucnl}
\def\universitydepartmentimagecfg {height=1.57cm}
\def\universitylocation {Monterrey, Nuevo Leon}

% INTEGRANTES, PROFESORES Y FECHAS
\def\authortable {
	\begin{tabular}{ll}
		\textbf{Alumno:}
		& \begin{tabular}[t]{l}
			 Martin Jonathan de la Cruz \\
		\end{tabular} \\
		Matricula: & UCNL \\

		\textit{Figura docente:}
		& \begin{tabular}[t]{l}
			Nombre \\ Apellido
		\end{tabular} \\
		& \\
		\multicolumn{2}{l}{Fecha de realizacion: \today} \\
		\multicolumn{2}{l}{\universitylocation}
	\end{tabular}
}

% IMPORTACION DEL TEMPLATE BASE
\input{template}

% FORMATO GENERAL
\usepackage{helvet}
\renewcommand{\familydefault}{\sfdefault}

% TABLAS Y PRODUCTOS VISUALES
\usepackage{array}
\usepackage{longtable}
\usepackage{pdflscape}
\usepackage{tikz}
\usetikzlibrary{
	arrows.meta,
	positioning,
	calc,
	matrix,
	fit,
	shapes.geometric,
	shadows.blur
}

% CITAS EN FORMATO AUTOR-ANIO, COMPATIBLE CON APA 7 EN EL CUERPO DEL TEXTO
% Usar:
% - Cita parentetica: \citep{clave}
% - Cita narrativa: \citet{clave}
% - Varias fuentes: \citep{clave1,clave2}
\setcitestyle{authoryear,open={(},close={)}}

% ==============================================================================
% AYUDAS DE REDACCION
% Quitar o comentar \pendiente cuando el documento final este terminado.
% ==============================================================================
\newcommand{\pendiente}[1]{\textcolor{red}{[PENDIENTE: #1]}}

% ==============================================================================
% MARCA DE AGUA EN PORTADA
% - Se aplica solo a la portada porque usa \AddToShipoutPictureBG*.
% - Para apagarla en alguna actividad: \def\coverwatermarkenabled {false}
% - Usar opacidad baja para que no compita con titulo, datos ni logotipo.
% ==============================================================================
\def\coverwatermarkenabled {true}
\def\coverwatermarkimage {img/departamentos/logo-nuevo-leon.png}
\def\coverwatermarkopacity {0.16}
\def\coverwatermarkwidth {0.72\paperwidth}
\def\coverwatermarkxshift {0cm}
\def\coverwatermarkyshift {-0.35cm}

\newcommand{\insertcoverwatermark}{%
	\ifthenelse{\equal{\coverwatermarkenabled}{true}}{%
		\AddToShipoutPictureBG*{%
			\begin{tikzpicture}[remember picture,overlay]
				\node[
					opacity=\coverwatermarkopacity,
					inner sep=0pt,
					xshift=\coverwatermarkxshift,
					yshift=\coverwatermarkyshift
				] at (current page.center) {%
					\includegraphics[width=\coverwatermarkwidth]{\coverwatermarkimage}%
				};
			\end{tikzpicture}%
		}%
	}{}%
}

\begin{document}

% PORTADA
\insertcoverwatermark
\templatePortrait

% CONFIGURACION DE PAGINA Y ENCABEZADOS
\templatePagecfg
\onehalfspacing

% ==============================================================================
% RESUMEN / ABSTRACT
% Debe ser breve. Formula recomendada:
% Tema + objetivo + tecnica/producto + enfoque personal.
% ==============================================================================
\begin{abstractd}
	\textbf{Objetivo:} Identificar los conceptos fundamentales del ciclo contable, la ecuacion contable, el control interno y los estados financieros mediante la resolucion de un cuestionario de opcion multiple.

	\textbf{Producto elaborado:} Cuestionario resuelto sobre libro diario, libro mayor, balanza de comprobacion, capital contable, costos y postulados basicos de las NIF.

	\textbf{Enfoque de analisis:} El cuestionario se responde desde una perspectiva funcional: cada reactivo se interpreta segun la operacion contable que explica, el control que exige y la forma en que impacta la calidad de la informacion financiera.

	\textbf{Nivel cognitivo:} Comprender y aplicar, porque los reactivos exigen reconocer conceptos contables y asociarlos con su uso correcto en registros, clasificacion de cuentas y elaboracion de estados financieros.
\end{abstractd}

% TABLA DE CONTENIDOS - INDICE
\templateIndex

% CONFIGURACIONES FINALES
\templateFinalcfg

% ======================= INICIO DEL DOCUMENTO =======================

% ==============================================================================
% SECCION 1: INTRODUCCION
% Apertura recomendada:
% - Iniciar con una afirmacion clara, no con una frase generica.
% - Traducir la consigna a un problema academico, juridico, etico, comunicativo
%   o profesional, segun la materia.
% - Explicar por que importa.
% - Cerrar anticipando el enfoque personal.
% ==============================================================================
\section{Introduccion}

La contabilidad no consiste solo en registrar operaciones: organiza la informacion economica de una entidad para convertirla en una base confiable de control, analisis y toma de decisiones. Cuando los registros no siguen una logica tecnica, la informacion financiera pierde utilidad y se debilita la capacidad de evaluar la situacion real de la empresa.

El problema que atiende esta actividad es distinguir con precision los conceptos que estructuran el ciclo contable y los estados financieros. No basta con reconocer terminos como activo, pasivo, capital, libro diario, libro mayor o balanza de comprobacion; es necesario entender la funcion que cada elemento cumple dentro del proceso de registro, clasificacion, verificacion y cierre contable.

Desde esta perspectiva, el cuestionario se resuelve con un criterio operativo: en cada reactivo se identifica primero el concepto contable involucrado, despues la funcion que cumple en el sistema de informacion y, por ultimo, la consecuencia que tiene sobre el control interno o la presentacion financiera.

% ==============================================================================
% SECCION 2: DESARROLLO / ANALISIS
% Adaptar los subtitulos segun la materia y la planeacion.
%
% Productos posibles:
% - Cuadro comparativo: criterios, evidencia, diferencia, valoracion.
% - Mapa conceptual: concepto central, conceptos secundarios, conectores claros.
% - Linea de tiempo: fecha/etapa, hecho, aporte, consecuencia.
% - Estudio de caso: situacion, elementos, problema, interpretacion, solucion.
% - Lista de verificacion: criterio observable, indicador, cumplimiento, mejora.
% - Glosario: concepto, definicion, funcion, ejemplo, fuente.
% - Reporte o ensayo: tesis, argumentos, evidencia, postura y cierre.
%
% Regla de parrafo:
% Entrada: concepto o problema.
% Proceso: explicacion, fuente y relacion.
% Salida: consecuencia, criterio o postura.
% ==============================================================================
\section{Desarrollo}

\subsection{Contexto y problema}

La Actividad 2 se ubica en la etapa introductoria de Contabilidad I, donde el estudiante debe consolidar los fundamentos del ciclo contable y comprender como se relacionan las cuentas con los estados financieros. El contexto del cuestionario exige pasar de la definicion aislada al uso correcto del concepto, ya que cada pregunta remite a una fase del proceso contable o a una decision de clasificacion.

El problema de analisis consiste en reconocer que la contabilidad funciona como un sistema ordenado: registra hechos economicos, los clasifica, verifica su equilibrio y finalmente los presenta en informes utiles. Por ello, responder correctamente no depende solo de memoria, sino de entender como se enlazan el diario, el mayor, la balanza, los ajustes y el cierre.

\subsection{Marco conceptual}

Los conceptos centrales del cuestionario son ciclo contable, ecuacion contable, capital contable, libro diario, libro mayor, balanza de comprobacion, balance general, estado de resultados y control interno. Cada uno cumple una funcion precisa dentro del sistema de informacion financiera. El libro diario registra cronologicamente las operaciones; el libro mayor concentra movimientos por cuenta; la balanza de comprobacion verifica el equilibrio entre cargos y abonos; y los estados financieros sintetizan la situacion economica de la entidad.

Tambien resultan relevantes los postulados basicos, como entidad economica y devengacion contable, porque delimitan cuando y como deben reconocerse los hechos economicos. A ello se suman las NIF y la clasificacion de cuentas en activo, pasivo y capital, que permiten mantener orden, comparabilidad y consistencia en la presentacion de la informacion. Sin ese marco conceptual, el registro contable se vuelve mecanico y pierde valor analitico.

\subsection{Nivel cognitivo aplicado}

El nivel cognitivo dominante es comprender y aplicar. Comprender implica distinguir que representa el activo, cual es la estructura del balance general o que establece el postulado de entidad economica; aplicar implica seleccionar la respuesta correcta a partir de la funcion real de cada cuenta, documento o procedimiento dentro del ciclo contable.

En consecuencia, el producto solicitado no se limita a repetir definiciones. Cada reactivo exige clasificar, interpretar y verificar la relacion entre operaciones, cuentas y estados financieros para sostener una respuesta contablemente correcta.

\subsection{Producto solicitado}

\textbf{Cuestionario resuelto sobre ciclo contable, control interno y estados financieros}

\begin{enumerate}
	\item \textbf{\textquestiondown Cada cuanto tiempo se realiza el ciclo contable completo?}
	\begin{itemize}
		\item a. Semanal
		\item b. Diario
		\item \textbf{c. Mensual, trimestral o anual segun la entidad}
		\item d. Cada cinco anos
	\end{itemize}
	\textbf{Respuesta correcta:} c. Mensual, trimestral o anual segun la entidad.

	\item \textbf{\textquestiondown Que parte de la ecuacion contable refleja la inversion de los propietarios?}
	\begin{itemize}
		\item \textbf{a. Capital contable}
		\item b. Gastos
		\item c. Activo
		\item d. Pasivo
	\end{itemize}
	\textbf{Respuesta correcta:} a. Capital contable.

	\item \textbf{\textquestiondown Que documento inicia formalmente el ciclo contable?}
	\begin{itemize}
		\item a. Libro mayor
		\item b. Balanza de comprobacion
		\item c. Estado de resultados
		\item \textbf{d. Libro diario}
	\end{itemize}
	\textbf{Respuesta correcta:} d. Libro diario.

	\item \textbf{\textquestiondown Que sucede cuando una empresa registra perdidas?}
	\begin{itemize}
		\item a. Aumenta el pasivo
		\item b. Aumenta el activo
		\item c. No afecta el patrimonio
		\item \textbf{d. Disminuye el capital}
	\end{itemize}
	\textbf{Respuesta correcta:} d. Disminuye el capital.

	\item \textbf{\textquestiondown Que contiene cada cuenta en el libro mayor?}
	\begin{itemize}
		\item a. La utilidad neta
		\item \textbf{b. Cargos, abonos y saldo}
		\item c. Unicamente abonos
		\item d. Solo el saldo final
	\end{itemize}
	\textbf{Respuesta correcta:} b. Cargos, abonos y saldo.

	\item \textbf{\textquestiondown Por que es importante el ciclo contable?}
	\begin{itemize}
		\item \textbf{a. Facilita la organizacion y presentacion de la informacion financiera}
		\item b. Controla las ventas
		\item c. Permite calcular impuestos
		\item d. Administra la nomina
	\end{itemize}
	\textbf{Respuesta correcta:} a. Facilita la organizacion y presentacion de la informacion financiera.

	\item \textbf{\textquestiondown Que tipo de cuentas aparecen en el libro mayor?}
	\begin{itemize}
		\item \textbf{a. Todas las cuentas contables}
		\item b. Solo cuentas de capital
		\item c. Solo cuentas de activo
		\item d. Unicamente cuentas de resultados
	\end{itemize}
	\textbf{Respuesta correcta:} a. Todas las cuentas contables.

	\item \textbf{\textquestiondown Que es el control interno en contabilidad?}
	\begin{itemize}
		\item a. Revision de impuestos
		\item b. Supervision gubernamental
		\item \textbf{c. Procedimientos para salvaguardar los recursos de la empresa}
		\item d. Auditoria externa
	\end{itemize}
	\textbf{Respuesta correcta:} c. Procedimientos para salvaguardar los recursos de la empresa.

	\item \textbf{\textquestiondown Como se clasifican las cuentas en contabilidad?}
	\begin{itemize}
		\item a. Activo, pasivo y resultados
		\item \textbf{b. Activo, pasivo y capital}
		\item c. Ingresos, egresos y ganancias
		\item d. Efectivo, inventario y deudas
	\end{itemize}
	\textbf{Respuesta correcta:} b. Activo, pasivo y capital.

	\item \textbf{\textquestiondown El capital contable representa:}
	\begin{itemize}
		\item a. Las obligaciones de la empresa
		\item b. Los ingresos acumulados
		\item \textbf{c. La inversion de los propietarios mas las utilidades retenidas}
		\item d. Los recursos disponibles
	\end{itemize}
	\textbf{Respuesta correcta:} c. La inversion de los propietarios mas las utilidades retenidas.

	\item \textbf{\textquestiondown Las NIF son aplicables a que tipo de entidades?}
	\begin{itemize}
		\item \textbf{a. A todas las entidades economicas}
		\item b. Solo empresas privadas
		\item c. Unicamente instituciones financieras
		\item d. Solo empresas publicas
	\end{itemize}
	\textbf{Respuesta correcta:} a. A todas las entidades economicas.

	\item \textbf{\textquestiondown Que principio establece que una entidad es independiente de sus propietarios?}
	\begin{itemize}
		\item \textbf{a. Entidad economica}
		\item b. Prudencia
		\item c. Consistencia
		\item d. Uniformidad
	\end{itemize}
	\textbf{Respuesta correcta:} a. Entidad economica.

	\item \textbf{\textquestiondown Cual es el objetivo principal de la balanza de comprobacion?}
	\begin{itemize}
		\item \textbf{a. Verificar el equilibrio contable}
		\item b. Calcular impuestos
		\item c. Preparar estados financieros
		\item d. Registrar compras
	\end{itemize}
	\textbf{Respuesta correcta:} a. Verificar el equilibrio contable.

	\item \textbf{\textquestiondown Cual es una caracteristica del control interno efectivo?}
	\begin{itemize}
		\item a. Registro manual de operaciones
		\item b. Falta de supervision
		\item \textbf{c. Separacion de funciones}
		\item d. Flexibilidad total en procesos
	\end{itemize}
	\textbf{Respuesta correcta:} c. Separacion de funciones.

	\item \textbf{\textquestiondown El balance general muestra:}
	\begin{itemize}
		\item a. La utilidad neta
		\item b. El flujo de efectivo mensual
		\item \textbf{c. La situacion financiera en una fecha determinada}
		\item d. Los ingresos y gastos
	\end{itemize}
	\textbf{Respuesta correcta:} c. La situacion financiera en una fecha determinada.

	\item \textbf{\textquestiondown Cual es la estructura basica del balance general?}
	\begin{itemize}
		\item a. Efectivo + Inventarios
		\item \textbf{b. Activo $=$ Pasivo + Capital}
		\item c. Ventas-Costos
		\item d. Ingresos - Gastos = Utilidad
	\end{itemize}
	\textbf{Respuesta correcta:} b. Activo $=$ Pasivo + Capital.

	\item \textbf{\textquestiondown Cual de las siguientes es una cuenta de activo?}
	\begin{itemize}
		\item a. Capital social
		\item \textbf{b. Caja}
		\item c. Prestamos bancarios
		\item d. Cuentas por pagar
	\end{itemize}
	\textbf{Respuesta correcta:} b. Caja.

	\item \textbf{\textquestiondown Cual es el efecto del asiento de cierre sobre la cuenta de resultados?}
	\begin{itemize}
		\item a. La elimina
		\item b. Disminuye el pasivo
		\item \textbf{c. Transfiere la utilidad o perdida al capital contable}
		\item d. Duplica los ingresos
	\end{itemize}
	\textbf{Respuesta correcta:} c. Transfiere la utilidad o perdida al capital contable.

	\item \textbf{\textquestiondown Que representa el "activo" en la ecuacion contable?}
	\begin{itemize}
		\item a. Las obligaciones financieras
		\item \textbf{b. Los recursos que posee la empresa}
		\item c. El capital invertido por los socios
		\item d. Las deudas de la empresa
	\end{itemize}
	\textbf{Respuesta correcta:} b. Los recursos que posee la empresa.

	\item \textbf{\textquestiondown Cual de los siguientes es un ejemplo de costo indirecto?}
	\begin{itemize}
		\item \textbf{a. Sueldo del personal administrativo}
		\item b. Mano de obra directa
		\item c. Materia prima
		\item d. Insumos directos
	\end{itemize}
	\textbf{Respuesta correcta:} a. Sueldo del personal administrativo.

	\item \textbf{\textquestiondown Que representa el costo primo?}
	\begin{itemize}
		\item a. Los costos fijos totales
		\item \textbf{b. La suma de materia prima y mano de obra directa}
		\item c. La suma de costos indirectos
		\item d. Los costos de distribucion
	\end{itemize}
	\textbf{Respuesta correcta:} b. La suma de materia prima y mano de obra directa.

	\item \textbf{\textquestiondown La contabilidad se define como:}
	\begin{itemize}
		\item \textbf{a. La tecnica que produce informacion financiera}
		\item b. La ciencia social que estudia el comportamiento humano
		\item c. Un metodo de control fiscal
		\item d. Un registro de facturas
	\end{itemize}
	\textbf{Respuesta correcta:} a. La tecnica que produce informacion financiera.

	\item \textbf{\textquestiondown Que ocurre si una cuenta se carga?}
	\begin{itemize}
		\item \textbf{a. Depende del tipo de cuenta}
		\item b. No afecta el saldo
		\item c. Siempre aumenta
		\item d. Siempre disminuye
	\end{itemize}
	\textbf{Respuesta correcta:} a. Depende del tipo de cuenta.

	\item \textbf{\textquestiondown Cual de las siguientes es una caracteristica de la informacion financiera segun las NIF?}
	\begin{itemize}
		\item a. Ser exclusiva para accionistas
		\item b. Ser compleja
		\item \textbf{c. Ser confiable y comparable}
		\item d. Ser confidencial
	\end{itemize}
	\textbf{Respuesta correcta:} c. Ser confiable y comparable.

	\item \textbf{\textquestiondown En que estado financiero se presentan los ingresos y gastos?}
	\begin{itemize}
		\item a. Balance general
		\item b. Estado de flujo de efectivo
		\item \textbf{c. Estado de resultados}
		\item d. Estado de cambios en el capital
	\end{itemize}
	\textbf{Respuesta correcta:} c. Estado de resultados.

	\item \textbf{\textquestiondown Que postulado establece que se deben registrar los hechos economicos en el momento en que ocurren?}
	\begin{itemize}
		\item a. Consistencia
		\item b. Entidad economica
		\item \textbf{c. Devengacion contable}
		\item d. Sustancia economica
	\end{itemize}
	\textbf{Respuesta correcta:} c. Devengacion contable.

	\item \textbf{\textquestiondown Cual es el otro nombre del balance general?}
	\begin{itemize}
		\item a. Estado patrimonial
		\item \textbf{b. Estado de situacion financiera}
		\item c. Estado de perdidas y ganancias
		\item d. Estado de liquidez
	\end{itemize}
	\textbf{Respuesta correcta:} b. Estado de situacion financiera.

	\item \textbf{\textquestiondown Que se elabora antes del cierre contable?}
	\begin{itemize}
		\item a. Estado de flujo de efectivo
		\item b. Diario general
		\item c. Estado de resultados
		\item \textbf{d. Balanza de comprobacion ajustada}
	\end{itemize}
	\textbf{Respuesta correcta:} d. Balanza de comprobacion ajustada.

	\item \textbf{\textquestiondown Que es un costo variable?}
	\begin{itemize}
		\item \textbf{a. El que depende del volumen de produccion}
		\item b. El que no cambia nunca
		\item c. El costo de impuestos
		\item d. El costo de administracion
	\end{itemize}
	\textbf{Respuesta correcta:} a. El que depende del volumen de produccion.

	\item \textbf{\textquestiondown Cual de los siguientes es un ejemplo de pasivo?}
	\begin{itemize}
		\item a. Inventario
		\item b. Cuentas por cobrar
		\item \textbf{c. Prestamo bancario}
		\item d. Caja
	\end{itemize}
	\textbf{Respuesta correcta:} c. Prestamo bancario.
\end{enumerate}

\subsection{Tecnica didactica aplicada}

En esta actividad se utilizo la tecnica de cuestionario estructurado, porque permite verificar de manera puntual si el estudiante distingue funciones, documentos y relaciones basicas del sistema contable. La secuencia de reactivos organiza el contenido por conceptos clave y favorece la identificacion rapida de procedimientos, cuentas y estados financieros, por lo que el producto no se limita a enlistar respuestas, sino que demuestra comprension aplicada del tema.

% ==============================================================================
% BLOQUE OPCIONAL: TABLA O CUADRO COMPARATIVO AMPLIO
% Usar cuando el producto necesite mucho ancho:
% - \clearpage para iniciar en pagina limpia.
% - \pagestyle{empty} para quitar encabezado y pie en el producto amplio.
% - landscape para aprovechar el ancho horizontal.
% - scriptsize, tabcolsep y arraystretch para controlar legibilidad.
% - Restaurar \pagestyle{fancy} al terminar.
%
% Si la tabla es corta, puede omitirse landscape y usarse longtable normal.
% ==============================================================================
% \clearpage
% \pagestyle{empty}
% \begin{landscape}
% \begingroup
% \scriptsize
% \setlength{\tabcolsep}{4pt}
% \renewcommand{\arraystretch}{1.35}
% \begin{longtable}{@{}p{0.18\linewidth}p{0.39\linewidth}p{0.39\linewidth}@{}}
% \caption{Titulo del cuadro comparativo}\label{tab:cuadro-comparativo}\\
% \toprule
% \textbf{Criterio} & \textbf{Elemento 1} & \textbf{Elemento 2} \\
% \midrule
% \endfirsthead
% \toprule
% \textbf{Criterio} & \textbf{Elemento 1} & \textbf{Elemento 2} \\
% \midrule
% \endhead
% Concepto & \pendiente{Contenido sintetico con cita si aplica.} & \pendiente{Contenido sintetico con cita si aplica.} \\
% Caracteristicas & \pendiente{Contenido.} & \pendiente{Contenido.} \\
% Funcion & \pendiente{Contenido.} & \pendiente{Contenido.} \\
% Valoracion critica & \pendiente{Contenido.} & \pendiente{Contenido.} \\
% \bottomrule
% \end{longtable}
% \endgroup
% \end{landscape}
% \pagestyle{fancy}

% ==============================================================================
% BLOQUE OPCIONAL: PRODUCTO VISUAL EN TIKZ
% Usar para mapa conceptual, mapa mental, esquema, flujo, proceso o diagrama.
% Cada flecha debe expresar relacion; no usar nodos sueltos sin conectores.
% ==============================================================================
% \begin{figure}[H]
% 	\centering
% 	\begin{tikzpicture}[
% 		node distance=1.3cm,
% 		box/.style={rectangle, rounded corners=2pt, draw=black!65, fill=black!4, align=center, minimum width=3.2cm, minimum height=0.9cm},
% 		arrow/.style={-{Latex[length=2.5mm]}, thick, draw=black!65}
% 	]
% 		\node[box] (centro) {Concepto central};
% 		\node[box, right=of centro] (relacion) {Concepto relacionado};
% 		\node[box, below=of relacion] (efecto) {Consecuencia};
%
% 		\draw[arrow] (centro) -- node[above]{implica} (relacion);
% 		\draw[arrow] (relacion) -- node[right]{produce} (efecto);
% 	\end{tikzpicture}
% 	\caption{Titulo del producto visual.}
% \end{figure}

\subsection{Analisis propio}

El principal hallazgo del cuestionario es que la contabilidad se sostiene en una logica de orden y verificacion continua. Los reactivos muestran que registrar operaciones, clasificarlas correctamente y trasladarlas a estados financieros no son tareas aisladas, sino partes de un mismo proceso de control. La interpretacion que se desprende es que la informacion financiera solo resulta util cuando se construye con consistencia tecnica.

Tambien se observa que conceptos como capital contable, control interno, devengacion y balanza de comprobacion no son definiciones decorativas, sino mecanismos que permiten mantener equilibrio, detectar errores y evaluar la posicion financiera de la empresa. La consecuencia de comprender esta articulacion es que la contabilidad puede utilizarse como herramienta de gestion y no solo como registro historico.

% ==============================================================================
% SECCION 3: POSTURA PERSONAL
% Debe mostrar comprension real. No repetir fuentes; interpretar.
% ==============================================================================
\section{Postura personal}

Considero que la contabilidad adquiere verdadero valor cuando se entiende como un sistema de control y decision, y no solamente como una tecnica para cumplir obligaciones formales. La razon es que cada cuenta, cada asiento y cada estado financiero resume hechos economicos que afectan patrimonio, liquidez y sostenibilidad. La consecuencia de asumir esta perspectiva es que el registro contable se vuelve una base confiable para dirigir mejor la empresa y anticipar riesgos.

% ==============================================================================
% SECCION 4: CONCLUSION
% Cierre recomendado:
% - Responder directamente a la actividad.
% - Indicar que se comprendio, no solo que se elaboro.
% - Mencionar criterio personal final.
% - No introducir informacion nueva.
% ==============================================================================
\section{Conclusion}

La Actividad 2 permitio resolver un cuestionario centrado en los elementos esenciales del ciclo contable, la ecuacion contable, el control interno y los estados financieros. A partir de los reactivos, se comprendio que el libro diario inicia el registro formal, que el libro mayor organiza movimientos por cuenta, que la balanza verifica el equilibrio contable y que el cierre transfiere el resultado al capital contable. En sintesis, el cuestionario confirma que la contabilidad no solo registra operaciones: estructura informacion confiable para controlar recursos, interpretar resultados y apoyar decisiones.

% ==============================================================================
% MANUAL DE TECNICAS DIDACTICAS
% Dejar incluido en la plantilla maestra. En una entrega final, puede comentarse
% esta linea si la docente solo solicita el producto de la semana.
% ==============================================================================
% Referencia opcional desactivada: archivo externo no disponible en este repositorio.

% BIBLIOGRAFIA
% Cambiar el nombre del archivo .bib segun la materia/actividad.
% Ejemplo: \bibliography{filosofia-del-derecho}
\bibliography{contabilidad-I}

% FIN DEL DOCUMENTO
\end{document}


